% ===== §2 A Critical Map of Wealth Theories =====
\section{A Critical Map of Wealth Theories}
\label{p15:sec:wealth_theories}

\noindent
Before we can develop the concept of Generative Relational Wealth, we must survey the landscape of existing wealth theories and ask, of each: what does it illuminate about wealth in intimate relations, and what does it systematically miss? The survey reveals a striking convergence: every major tradition in the theory of wealth shares, at the level of its fundamental ontology, the assumption that the primary subject of wealth is the individual. The concept of GRW requires a different starting point: a relational ontology in which wealth belongs primarily to the relational field rather than to any individual participant.

\subsection{Mercantilism: Wealth as Accumulation and Zero-Sum Exchange}
\label{p15:subsec:mercantilism}

\noindent
Mercantilism holds that wealth consists in the accumulation of precious metals, and that trade is a zero-sum competition: one party's gain is necessarily another's loss \parencite{marx1976}. Three features of mercantilist wealth are directly relevant to intimate relations: it is \emb{material and finite} (discrete, countable objects that can be possessed and transferred); \emb{zero-sum} (the total stock is fixed); and \emb{externally verifiable} (measurable by any competent observer, independently of the parties' feelings).

When mercantilist logic enters intimate relations---most visibly in bride-price and in the use of gold as the medium of relational contracts---it imports all three features. The bride-price treats relational value as a finite, measurable, transferable quantity. The exchange is zero-sum: accounts balance when the agreed gold changes hands. What mercantilist logic cannot see is the generation of relational value within the intimate relation itself---the accumulation, through co-evolutionary activity, of something that was not there before and that belongs to neither party individually. For mercantilism, wealth precedes the relation and is transferred through it; wealth generated by the relation, belonging to the relational field itself, is outside the conceptual framework.

\begin{claim}[The mercantilist blind spot]
Mercantilist wealth logic treats relational value as a finite, transferable, externally verifiable quantity that precedes and is exchanged through the relation. It is constitutively unable to recognise the generation of relational wealth through co-evolutionary activity---wealth that belongs to the \emph{between} rather than to either party, and that is augmented rather than depleted by use.
\end{claim}

\subsection{Classical Political Economy: Labour, Division, and Relational Surplus}
\label{p15:subsec:classical}

\noindent
Adam Smith's innovation was to locate the source of wealth not in the stock of precious metals but in \emb{labour}: the productive activity through which human beings transform the natural world \parencite{sen1999}. This is a genuine advance---wealth is generated, not merely transferred. But classical political economy remains limited by its focus on material production: the labour of generating relational value---the emotional labour, the care work, the attentional investment that builds a shared life---is invisible in this framework, because it does not produce material goods and therefore does not count as productive in the technical sense.

Ricardo's concept of \emb{differential advantage} will prove relevant when we consider the distribution of relational surplus: who captures the value generated by co-evolutionary activity in the intimate relational field, and whether that distribution is just.

\subsection{The Marxist Tradition: Commodity Fetishism, Alienation, and Social Reproduction}
\label{p15:subsec:marx}

\noindent
Marx's analysis of \emb{commodity fetishism} applies with striking directness to traditional wealth forms in intimate relations \parencite{marx1976}. The bride-price is a commodity fetishism of relational value: the capacity to participate in a shared life is mystified as a property of the gold exchanged. The social relations that produce and distribute this value are obscured behind the appearance of a simple exchange of equivalents.

Marx's concept of \emb{alienation} is equally relevant: the relational value generated through a person's emotional labour and care work is systematically appropriated by the dominant party, leaving the generating party estranged from the value they have produced.

The most important Marxist contribution comes from Silvia Federici's account of \emb{social reproduction} \parencite{federici2004}. Federici argues that capitalism depends on a vast domain of unpaid reproductive labour---emotional support, care work, household maintenance---performed predominantly by women and systematically excluded from the category of productive labour. This reproductive labour is not merely a precondition for productive labour; it is itself value production, and its non-recognition is both economic and epistemological injustice. The GRW framework provides the ontological grounding for Federici's analysis: what is being appropriated when emotional labour is exploited is not merely a service but a genuine form of wealth---irreducible to any commodity category, generated in the relational field, and belonging to the \emph{between}.

\begin{claim}[The Marxist contribution and its limits]
The Marxist tradition contributes two essential concepts: commodity fetishism (the mystification of relational value as a property of exchanged objects) and social reproduction (the systematic non-recognition of emotional and care labour as wealth production). Both converge with GRW theory. But the Marxist tradition remains limited by its individualist conception of the producing subject: it does not yet have the ontological resources to articulate wealth that belongs to the relational field itself.
\end{claim}

\subsection{The Capabilities Approach: Wealth as Freedom}
\label{p15:subsec:capabilities}

\noindent
Amartya Sen's capabilities approach reconceives what wealth is for \parencite{sen1999}. The goal of economic development is not the accumulation of material goods but the expansion of human capabilities---the real freedoms people have to lead lives they have reason to value. This resonates deeply with the GRW framework: the question to ask of any wealth form in intimate relations is whether it expands or constrains the relational capabilities of the parties involved.

More fundamentally, the capabilities approach points toward a reconception of relational wealth itself: not a stock of material goods but a set of capacities for relational flourishing---the capacity for joint attention, for timely sharing, for co-evolutionary receptivity, for non-possessive participation in the relational field's generative activity. These relational capabilities are the human substrate of GRW.

The capabilities approach remains limited, however, by its individualist framework: capabilities are conceived as properties of individual persons. The GRW framework requires a relational extension: capabilities that belong to the relational field itself, jointly held and jointly developed, whose expansion constitutes relational flourishing irreducible to the sum of individual capability expansions.

\subsection{The Gift Economy: Wealth as Flow and Social Bond}
\label{p15:subsec:gift}

\noindent
Marcel Mauss's analysis of the gift reveals a conception of wealth fundamentally different from commodity logic \parencite{mauss1990}. In gift economies, wealth is not accumulated but circulated: it flows through the social body, creating and sustaining the relations that constitute community life. Like GRW, gift wealth is essentially relational---its value is constituted by the social relations it creates, not by any intrinsic property of the material object. And like GRW, gift wealth cannot simply be accumulated: the gift that is hoarded loses its relational character.

But GRW differs from gift wealth in a crucial respect: the gift can be given, received, and reciprocated---it is, at least in principle, transferable between individual parties. GRW cannot be given; it can only be generated together. The happiness jar cannot be handed from one person to another; the shared attractor landscape of a life built together cannot be transferred to a new relationship. GRW is more radically relational than the gift: it does not just flow through relations but is constituted by them, making individual ownership or transfer conceptually impossible.

\subsection{Generative Justice: Wealth as Value Returned to Its Creators}
\label{p15:subsec:generative_justice}

\noindent
Ron Eglash's generative justice holds that value should be returned to those who generate it---that the extraction of value without return is injustice, and that just economic arrangements support and enhance the generative capacity of communities \parencite{eglash2016}. This applies directly to intimate relations: when emotional labour generates relational value that is appropriated by the dominant party rather than returned to the relational field, this is a generative justice failure.

But GRW extends this framework crucially. For Eglash, the generator of value is ultimately an individual or community. GRW holds that the primary generator of relational value is the relational field itself: the co-evolutionary dynamics produce wealth that belongs to the field, not to either individual participant. The generative justice question in intimate relations is therefore not just ``whose labour produced this?'' but ``what conditions allow the relational field itself to be a generative system, and are those conditions being sustained?''

\subsection{Non-Western Wealth Traditions}
\label{p15:subsec:nonwestern}

\noindent
The non-Western traditions of wealth offer important correctives to the individualist presuppositions of the Western tradition.

\emb{Ubuntu philosophy} (\emph{umuntu ngumuntu ngabantu}---a person is a person through other persons) holds that wealth, like personhood itself, is constituted through relations \parencite{mbiti1970,metz2011}. The isolated individual who accumulates wealth at the expense of community relations is not wealthy but impoverished: they have gained material goods at the cost of the relational substance that constitutes full personhood.

\emb{Daoist and Buddhist economics} converge on the insight that true wealth consists in non-possessive sufficiency rather than accumulation: \zh{知足者富}---those who know sufficiency are wealthy \parencite{laozi1963}. The non-possessive orientation toward relational wealth that the GRW framework recommends is the Daoist and Buddhist economic orientation applied to intimate life.

\emb{Confucian relational ethics} holds that true flourishing is inseparable from the health of the relational bonds that constitute the social fabric---an implicit theory of relational wealth that anticipates the GRW framework \parencite{tu1985}.

\begin{claim}[The convergence of non-Western traditions]
The non-Western wealth traditions converge on an insight the Western tradition has systematically suppressed: the deepest form of wealth is relational rather than material, generated through the cultivation of genuine relational bonds. GRW theory is, in part, the formal articulation of this ancient intuition within a contemporary philosophical and scientific framework.
\end{claim}

\subsection{The Shared Presupposition and Its Overcoming}
\label{p15:subsec:sharedpresupposition}

\noindent
Despite their profound differences, the major traditions in wealth theory share a structural feature that limits their application to intimate relations: they all conceive wealth primarily as a property of individuals---something that individuals have, produce, exchange, or are deprived of. Even the traditions most sensitive to the relational dimensions of wealth do not fully achieve the ontological shift that GRW requires: the relocation of the primary subject of wealth from the individual to the relational field.

This shared presupposition is not arbitrary. The most socially organised forms of wealth---the gold in the vault, the commodity in the marketplace, the capability of the individual person---are all, at least in principle, attributable to particular individuals or institutions. Relational wealth---wealth that belongs to the \emph{between}---is systematically invisible to frameworks that begin with the individual, because it cannot be attributed to any individual without distortion.

The next section examines how wealth forms organised by individualist logic enter the intimate relational field and affect its dynamics---and specifically, how three mechanisms of individualisation destroy the conditions for GRW generation.
